\chapter{Introduction}

This chapter describes the lexicon and syntax of the language.
Grammar rules\footnotemark[1] will be shown using a dialect of eBNF.

\footnotetext[1]{
	Only a subset of the rules will be shown, and they may not exactly match
	the real grammar. For the actual rules consult the appropriate grammar
	files.
}

\begin{listing}
	\label{lst:ebnf_dialect}
	\caption{Example of a definition of \code{some\_rule} using the eBNF dialect}
	\begin{minted}{antlr}
some_rule
	: 'some_token' some_label=some_other_rule
	| zero_or_one_of_this?
	| zero_or_more_of_this*
	| one_or_more_of_this+
	| ('some_token' | 'some_other_token' | some_other_rule)
	;
	\end{minted}
\end{listing}

\paragraph{Whitespaces} Whitespaces (e.g. spaces, tabs, newlines) are
determinant \textbf{only} during the lexical parsing, meaning that \code{!=} is
different than \code{! =}, where the former gets parsed as a single token while
the latter gets parsed as two different tokens; in any other case, whitepsaces
are ignored.

